\documentclass[12pt,a4paper]{article}
\usepackage[UTF8]{ctex}
\usepackage[margin=2.5cm]{geometry}
\usepackage{graphicx}
\usepackage{booktabs}
\usepackage{hyperref}
\usepackage{natbib}
\usepackage{amsmath}
\usepackage{setspace}
\usepackage{enumitem}

\bibliographystyle{apalike}
\setstretch{1.5}

\title{\textbf{人工智能在乳腺癌诊断中的研究进展综述}}
\author{MedgeClaw 学术写作系统}
\date{2026年3月}

\begin{document}
\maketitle

\begin{abstract}
乳腺癌是全球女性最常见的恶性肿瘤，早期精准诊断对改善患者预后至关重要。近年来，以深度学习为代表的人工智能（AI）技术在乳腺癌影像诊断领域取得了突破性进展。本文系统综述了AI技术在乳腺X线摄影、超声、磁共振成像（MRI）及组织病理学图像分析中的最新研究成果，涵盖了卷积神经网络（CNN）、Vision Transformer（ViT）等主流架构的应用，以及商业AI系统的临床验证结果。同时，本文讨论了多模态融合、可解释性AI（XAI）等前沿方向，并深入分析了数据隐私保护、模型泛化能力、临床落地障碍等关键挑战。通过对近年来高质量文献的梳理与总结，本文旨在为研究人员和临床医生提供该领域的全面参考，并展望AI辅助乳腺癌诊断的未来发展趋势。

\noindent\textbf{关键词：} 人工智能；深度学习；乳腺癌；医学影像诊断；计算机辅助检测
\end{abstract}

\section{引言}

乳腺癌是全球范围内对女性健康威胁最大的恶性肿瘤之一。根据国际癌症研究机构（IARC）发布的GLOBOCAN 2022数据，全球每年新增乳腺癌病例约230万例，占所有女性癌症新发病例的近四分之一，已超越肺癌成为全球发病率最高的癌种\citep{bray2024global}。在中国，乳腺癌同样呈现高发态势，且发病年龄有年轻化的趋势。早期诊断是降低乳腺癌死亡率的关键策略，研究表明，早期（I期）乳腺癌患者的五年生存率可超过90\%，而晚期患者的预后则显著恶化\citep{sung2021global}。

传统的乳腺癌筛查和诊断高度依赖于影像学检查及病理学评估，包括乳腺X线摄影（mammography）、超声成像、磁共振成像（MRI）以及组织病理学切片分析。然而，这些方法在实践中面临诸多挑战：乳腺X线摄影的假阳性率和假阴性率仍然偏高，致密型乳腺的检出率尤为不足；病理学诊断高度依赖专家经验，观察者间差异显著；此外，影像学医生的工作负荷日益加重，也制约了诊断效率和一致性。

人工智能（AI），特别是深度学习（deep learning）技术的快速发展，为突破上述瓶颈提供了新的路径。自2012年AlexNet在ImageNet竞赛中取得突破性成果以来\citep{krizhevsky2012imagenet}，深度学习在计算机视觉领域经历了爆发式增长\citep{lecun2015deep}。卷积神经网络（CNN）架构的不断演进——从ResNet\citep{he2016deep}到DenseNet\citep{huang2017densely}、EfficientNet\citep{tan2019efficientnet}，再到Vision Transformer\citep{dosovitskiy2021image}——为医学影像分析提供了越来越强大的特征提取和模式识别能力。在医学影像领域，\citet{litjens2017survey}的里程碑式综述系统总结了深度学习在医学图像分析中的广泛应用，奠定了该领域的研究基础。

近年来，AI在乳腺癌诊断中的应用研究呈指数级增长。多项大规模研究表明，AI系统在乳腺X线摄影判读中的诊断性能已接近甚至超过放射科医生\citep{mckinney2020international, rodriguez2019standalone}。在组织病理学领域，深度学习算法在检测淋巴结转移等任务上也展现出令人瞩目的精度\citep{bejnordi2017diagnostic}。这些突破性成果推动了AI辅助诊断系统从实验室研究向临床转化的进程。

本文旨在系统综述AI技术在乳腺癌诊断各个环节中的最新研究进展。文章将分别阐述AI在乳腺X线摄影、超声成像、磁共振成像和组织病理学图像分析四大模态中的应用现状，随后探讨多模态融合与可解释性AI的前沿研究方向，最后深入分析该领域面临的关键挑战并展望未来发展趋势。

\section{AI在乳腺X线摄影中的应用}

乳腺X线摄影是乳腺癌筛查的"金标准"，也是AI技术应用最为广泛和成熟的乳腺影像模态。深度学习模型在乳腺X线摄影中的应用主要集中在病灶检测与分类、癌症风险评估以及辅助筛查三个方面。

\subsection{深度学习架构与检测性能}

早期基于CNN的乳腺X线摄影分析系统主要采用经典的分类架构进行迁移学习。\citet{shen2019deep}提出了一种基于端到端深度学习的乳腺癌检测框架，在全视野数字乳腺X线摄影（FFDM）上实现了0.895的AUC，证明了深度学习直接处理全分辨率乳腺X线图像的可行性。随后，\citet{wu2020deep}开发的多视图深度神经网络在NYU数据集上进一步提升了检测性能，其系统在辅助放射科医生判读时显著降低了假阴性率。

具有里程碑意义的是\citet{mckinney2020international}在\textit{Nature}上发表的研究，该团队来自Google Health，基于大规模英美两国筛查数据开发的AI系统，在独立测试集上的AUC达到0.889（英国数据集），相比放射科医生，假阳性率降低了5.7\%，假阴性率降低了9.4\%。这一结果引发了学界对AI取代或辅助放射科医生进行乳腺癌筛查的广泛讨论。\citet{rodriguez2019standalone}的多中心研究进一步验证了这一趋势，其开发的AI系统在与101名放射科医生的对照中，诊断性能达到了放射科医生的中位水平。

近年来，基于Transformer架构的模型也开始在乳腺X线摄影分析中崭露头角。Vision Transformer（ViT）\citep{dosovitskiy2021image}通过自注意力机制捕获图像的全局上下文信息，在处理大尺寸医学图像时展现出独特优势。\citet{wang2024mammography}的综述详细总结了从传统CNN到Transformer的架构演进对乳腺X线摄影检测性能的影响。EfficientNet\citep{tan2019efficientnet}等高效架构则在保持高精度的同时显著降低了计算开销，使得AI系统在资源受限的临床环境中的部署成为可能。

\subsection{商业AI系统与临床试验}

随着AI技术在乳腺X线摄影中的不断成熟，多款商业化AI辅助诊断系统已获得FDA或CE认证并进入临床应用，包括Lunit INSIGHT Mammography、Transpara（ScreenPoint Medical）、CureMetrix cmAssist等。这些商业系统通常基于深度CNN架构，经过大规模标注数据训练，能够自动检测乳腺X线摄影中的可疑病灶并提供恶性概率评分。

在临床验证方面，2023年发表在\textit{The Lancet Oncology}上的MASAI试验是迄今为止规模最大的前瞻性随机对照研究之一\citep{lang2023artificial}。该试验在瑞典纳入了超过8万名筛查女性，比较了AI辅助单阅片方案与传统双阅片方案的筛查效能。结果表明，AI辅助方案的癌症检出率与传统方案相当（AI组6.7‰ vs 对照组5.0‰），同时将放射科医生的工作量降低了44.3\%。同年，\citet{dembrower2023artificial}在\textit{The Lancet Digital Health}上报告的另一项瑞典前瞻性研究也得出了类似结论，AI辅助筛查在不降低检出率的前提下显著提高了工作效率。

\citet{kim2020changes}的多读者回顾性研究从另一个角度验证了AI的临床价值：在AI辅助下，放射科医生的癌症检出率提升了约10\%，同时假阳性召回率有所下降。\citet{freeman2021use}发表在\textit{BMJ}上的系统综述全面分析了AI在乳腺癌筛查项目中的测试精度，指出虽然现有AI系统表现出色，但大多数研究仍存在方法学局限性，前瞻性临床验证的证据仍然不足。\citet{diaz2024artificial}进一步讨论了AI在数字乳腺X线摄影（DM）和数字乳腺断层合成（DBT）中的技术挑战与前景。

\subsection{风险评估与个性化筛查}

AI在乳腺癌领域的应用已超越了单纯的病灶检测，延伸到了癌症风险预测和个性化筛查策略的制定。\citet{yala2019deep}开发的基于深度学习的风险预测模型，利用乳腺X线摄影图像中的潜在影像组学特征（而非仅依赖传统临床风险因子），在预测未来5年乳腺癌风险方面显著优于Tyrer-Cuzick等传统模型。该团队后续的工作\citep{yala2022optimizing}进一步在国际多中心数据上验证了模型的泛化能力和鲁棒性。\citet{lowry2024artificial}和\citet{siddique2023tomography}分别从不同角度综述了AI在乳腺癌风险评估中的最新进展和应用前景。这些风险预测模型有望推动乳腺癌筛查从"一刀切"的群体策略向基于个体风险分层的精准筛查范式转变\citep{dembrower2020effect}。

\section{AI在乳腺超声诊断中的应用}

乳腺超声是乳腺X线摄影的重要补充，在致密型乳腺的评估中具有独特优势。AI在乳腺超声领域的应用主要聚焦于病灶检测、良恶性鉴别以及BI-RADS分类。

与乳腺X线摄影相比，超声图像具有更大的操作者依赖性和图像变异性，这对AI模型的鲁棒性提出了更高要求。早期研究中，\citet{han2022breast}提出了一种基于深度学习的乳腺超声病灶分类框架，利用GoogLeNet和ResNet等预训练模型进行迁移学习，在区分良恶性肿块方面取得了显著成效。\citet{byra2020breast}则探索了颜色空间转换结合迁移学习的策略，通过将灰度超声图像转换为伪彩色图像来充分利用在自然图像上预训练的CNN模型，该方法在独立测试集上的AUC达到0.936。\citet{fujioka2020distinction}的研究进一步验证了CNN在乳腺超声良恶性鉴别中的有效性。

近年来的研究趋势表明，融合多尺度特征和注意力机制的网络架构在乳腺超声分析中表现出更优的性能。\citet{zhang2023recent}在其综述中详细总结了AI在乳腺超声图像增强、分割和诊断方面的最新进展。\citet{alkarawi2024review}全面回顾了AI在乳腺成像各模态中的应用现状，指出超声AI系统在标准化评估流程方面仍面临挑战，主要原因在于缺乏类似于DDSM或VinDr等乳腺X线摄影领域的大规模公开标注数据集。尽管如此，自动化乳腺超声（ABUS）技术的普及为AI提供了更加标准化的图像采集方式，有望加速该领域的发展。

\section{AI在乳腺MRI中的应用}

乳腺MRI以其极高的灵敏度被广泛用于高危人群筛查、新辅助化疗疗效评估及手术范围规划等场景。然而，MRI检查的高特异性不足问题和图像解读的复杂性为AI的介入提供了广阔空间。

\citet{abdullah2025deep}在2025年发表的系统综述与Meta分析中，全面评估了深度学习模型在乳腺MRI诊断中的表现。该研究纳入了多项高质量研究，结果显示基于深度学习的模型在乳腺MRI病灶分类任务中的汇总AUC达到0.87--0.91，灵敏度和特异度均表现出色。\citet{adam2023deep}的文献综述从技术角度分析了不同深度学习架构在乳腺MRI中的应用，涵盖了从动态对比增强MRI（DCE-MRI）到扩散加权成像（DWI）的多种序列。\citet{truhn2019radiomic}的研究则比较了传统影像组学方法与深度学习在多参数乳腺MRI分类中的效能，结果表明深度学习方法在无需人工特征工程的情况下实现了更优的分类性能。

在功能性MRI方面，AI模型已被应用于预测新辅助化疗的病理完全缓解（pCR），这对于治疗决策的制定具有重要的临床意义。\citet{qi2024radiomics}在其关于乳腺癌影像组学的综述中，详细讨论了AI辅助的MRI影像组学在分子亚型识别、疗效预测和预后评估中的多维度应用。这些研究表明，AI不仅能够提升MRI的诊断准确性，还能从影像数据中挖掘出超越人类视觉感知的深层信息。

\section{AI在组织病理学图像分析中的应用}

组织病理学检查是乳腺癌确诊的"金标准"，而数字病理学的发展和全切片图像（Whole Slide Image, WSI）扫描技术的成熟为AI在该领域的应用奠定了基础。

\subsection{全切片图像分析与弱监督学习}

WSI的超高分辨率（通常为$10^9$级像素）对传统深度学习方法构成了严峻挑战。\citet{bejnordi2017diagnostic}在CAMELYON16挑战赛中的里程碑式研究表明，深度学习算法在检测乳腺癌前哨淋巴结转移方面可达到接近甚至超过病理医生的诊断精度。这项发表于\textit{JAMA}的研究极大地推动了AI在计算病理学中的应用。

为解决WSI的海量数据与有限标注之间的矛盾，多实例学习（Multiple Instance Learning, MIL）成为计算病理学中的主流范式。\citet{campanella2019clinical}在\textit{Nature Medicine}上发表的研究展示了弱监督深度学习在临床级WSI分析中的潜力，仅使用切片级标签即可实现前列腺癌、基底细胞癌和乳腺癌转移的高精度检测，AUC超过0.98。\citet{lu2021data}提出的CLAM（Clustering-constrained Attention Multiple Instance Learning）框架进一步推进了数据高效的弱监督病理图像分析，该方法通过注意力机制自动识别诊断相关的关键区域，在多种癌症类型上均展现出优异性能。

\subsection{Transformer与基础模型}

随着Vision Transformer\citep{dosovitskiy2021image}和自监督学习的兴起，计算病理学领域迎来了新一轮技术革新。\citet{chen2022scaling}提出的HIPT（Hierarchical Image Pyramid Transformer）通过分层自监督预训练策略将ViT扩展到千兆像素级的病理图像，在生存预测和分子亚型分类等任务上取得了显著突破。大规模病理基础模型（Foundation Models）如UNI、Virchow、CONCH等相继涌现，这些模型在海量未标注病理图像上进行预训练，展现出强大的迁移学习能力。

\citet{jiang2024deep}的综述系统总结了深度学习在乳腺癌组织病理学图像中的诊断、治疗和预后三个维度的应用，指出当前研究正从单纯的肿瘤检测向分子分型预测、微环境分析和疗效评估等更深层次的任务拓展。\citet{priya2024deep}的综述则从方法学角度全面梳理了深度学习在乳腺癌组织病理学图像检测中的各种技术路线，包括监督学习、半监督学习和无监督学习等范式。

\subsection{生物标志物预测}

AI在组织病理学中的另一个重要应用方向是从常规HE染色切片中预测分子生物标志物。HER2状态、Ki-67增殖指数、雌激素受体（ER）和孕激素受体（PR）表达状态等关键生物标志物的检测传统上需要免疫组化（IHC）或荧光原位杂交（FISH）等额外检测手段。深度学习模型已被证明可以直接从HE切片的形态学特征中预测这些标志物的表达状态，从而有望简化临床工作流程并降低检测成本\citep{jiang2024deep, arunkumar2023review}。

\section{多模态融合与可解释性AI}

\subsection{多模态数据融合}

单一影像模态的诊断信息存在局限性，融合多种影像模态（如乳腺X线摄影+超声+MRI）以及临床信息（年龄、家族史、基因检测结果等）的多模态AI系统正成为研究热点。\citet{zhang2023recent}指出，多模态融合策略主要包括早期融合（特征级拼接）、晚期融合（决策级融合）和中间融合（注意力引导的跨模态交互）三种范式。\citet{qi2024radiomics}的综述强调了影像组学在多模态数据整合中的桥梁作用，通过提取标准化的定量特征实现不同模态之间的互补。\citet{alkarawi2024review}进一步讨论了多模态AI在乳腺癌诊断中的实际应用案例和面临的数据对齐、标准化等技术挑战。

\subsection{可解释性AI}

深度学习模型常被视为"黑箱"，其决策过程的不透明性严重制约了AI系统在临床实践中的接受度和信任度。可解释性AI（Explainable AI, XAI）旨在使模型的推理过程对人类可理解，这对于医学AI的临床部署至关重要。

\citet{ghasemi2024explainable}在2024年发表的系统综述全面梳理了XAI在乳腺癌检测和风险预测中的应用现状。主流的XAI方法包括基于梯度的可视化技术（如Grad-CAM）\citep{selvaraju2017grad}、模型无关的解释方法（如LIME）\citep{ribeiro2016should}，以及基于注意力机制的内在可解释性设计。在乳腺影像分析中，Grad-CAM等方法通过生成热力图直观地显示模型关注的图像区域，帮助放射科医生理解和验证AI的判断依据。\citet{nasser2023deep}的系统综述也强调了XAI在推动深度学习乳腺癌诊断系统临床落地过程中的关键作用。

\section{挑战与展望}

\subsection{数据隐私与联邦学习}

大规模医学影像数据的收集和共享面临严格的隐私保护法规（如GDPR、HIPAA）约束。联邦学习（Federated Learning）\citep{mcmahan2017communication}作为一种分布式机器学习范式，允许多个医疗机构在不共享原始数据的前提下协作训练AI模型，为解决数据孤岛问题提供了重要的技术路径\citep{rieke2020future}。目前，多个国际合作项目已在乳腺癌诊断场景中验证了联邦学习的有效性，但通信效率、数据异质性和模型安全性等问题仍需进一步研究。

\subsection{模型泛化与外部验证}

AI模型的泛化能力是制约其临床推广的核心问题之一。由于不同医疗机构的设备型号、成像协议、患者群体构成等存在显著差异，在单一数据集上训练的模型往往难以在外部数据上保持同等性能。\citet{mckinney2020international}的研究之所以具有里程碑意义，部分原因在于其在英国和美国两个独立数据集上进行了跨国验证。\citet{freeman2021use}在系统综述中特别指出，现有大多数AI乳腺癌诊断研究缺乏充分的外部验证，方法学偏倚风险较高。未来研究应更加注重多中心、前瞻性临床试验的设计和实施。

\subsection{临床落地障碍}

尽管AI在乳腺癌诊断中的技术性能已相当出色，从研究到临床常规应用仍存在多重障碍。首先，监管审批流程复杂且各国标准不一，AI医疗器械的认证周期较长。其次，AI系统与现有临床工作流程的整合需要考虑IT基础设施、数据互操作性以及临床医生的培训和适应等实际问题。再者，AI辅助诊断的法律责任归属、医疗保险覆盖等制度性问题尚未明确。\citet{diaz2024artificial}深入讨论了AI乳腺癌检测技术从研发到临床部署面临的多维度挑战。\citet{almuhaisen2024ai}的综述也从实际临床角度分析了AI在乳腺X线摄影中应用的障碍与机遇。

\subsection{未来发展方向}

展望未来，AI在乳腺癌诊断中的发展将呈现以下主要趋势。第一，大规模病理和影像基础模型（Foundation Models）将推动计算病理学和放射影像学的范式变革，通过自监督预训练和少样本微调降低对大规模标注数据的依赖\citep{chen2022scaling}。第二，多组学数据（基因组、转录组、蛋白组）与影像数据的深度融合将实现从"影像诊断"到"精准医学"的跨越。第三，生成式AI（如扩散模型）在医学图像合成、数据增强方面的应用将有效缓解稀有病例数据不足的问题。第四，联邦学习与隐私计算技术的成熟将推动跨机构、跨国界的大规模协作研究\citep{rieke2020future}。第五，人机协作模式的优化——AI作为放射科医生的"智能助手"而非替代者——将是临床落地的主要路径\citep{lang2023artificial, dembrower2023artificial}。

\section{结论}

人工智能技术，特别是深度学习，已在乳腺癌诊断的多个环节展现出巨大的应用潜力。在乳腺X线摄影领域，AI系统的诊断性能已接近或超过放射科医生水平，大规模前瞻性临床试验进一步验证了其在真实筛查场景中的有效性和安全性。在超声和MRI领域，AI同样展现出提升诊断效率和准确性的能力。在组织病理学领域，弱监督学习和基础模型的发展正在推动计算病理学的快速进步。与此同时，多模态融合和可解释性AI的研究为构建更加全面、可信的诊断决策支持系统提供了技术基础。

然而，数据隐私保护、模型泛化验证、临床工作流程整合以及监管审批等挑战仍然存在。未来的研究应更加注重多中心前瞻性验证、联邦学习等隐私保护技术的应用，以及人机协作模式的优化设计。随着技术的不断成熟和临床证据的持续积累，AI辅助乳腺癌诊断有望在不远的将来成为临床常规实践的重要组成部分，为提升乳腺癌的早期检出率和患者预后做出实质性贡献。

\bibliography{../references/references}

\end{document}
